\documentclass{article}
\usepackage{amsmath,amssymb,amsthm}
\usepackage{algorithm,algorithmic}
\usepackage{bm}
\usepackage{hyperref}
\usepackage{enumitem}
\newtheorem{theorem}{Theorem}
\newtheorem{lemma}{Lemma}
\newtheorem{assumption}{Assumption}
\newtheorem{corollary}{Corollary}
\newcommand{\E}{\mathbb{E}}
\newcommand{\Var}{\operatorname{Var}}
\newcommand{\norm}[1]{\left\|#1\right\|}
\newcommand{\ip}[2]{\langle #1,#2\rangle}
\newcommand{\Id}{\mathbf I}
\begin{document}

\section*{Algorithm Name}
\textbf{Stochastic Subsampled Newton (SSN)}  

The method uses a stochastic estimate of the Hessian matrix obtained by sampling a mini‑batch of data at each iteration.  

\section*{Mathematical Formulation}
Let $f:\mathbb{R}^d\to\mathbb{R}$ be the empirical risk 
\[
f(x)=\frac1n\sum_{i=1}^n \ell_i(x),
\]
where each component loss $\ell_i$ is twice continuously differentiable.  
At iteration $t$ we draw a random subset $S_t\subset\{1,\dots,n\}$ of size $|S_t|=b$ and define  

\[
\begin{aligned}
g_t &:= \nabla f_{S_t}(x_t)=\frac1b\sum_{i\in S_t}\nabla \ell_i(x_t),\\[4pt]
H_t &:= \nabla^2 f_{S_t}(x_t)+\lambda \Id
      =\frac1b\sum_{i\in S_t}\nabla^2 \ell_i(x_t)+\lambda \Id,
\end{aligned}
\]
where $\lambda>0$ is a damping parameter guaranteeing $H_t\succ0$.  
The update rule is  

\[
\boxed{\,x_{t+1}=x_t-\alpha H_t^{-1}g_t\,},
\]
with step‑size (learning rate) $\alpha\in(0,1]$.

\section*{Pseudocode}
\begin{algorithm}[H]
\caption{Stochastic Subsampled Newton (SSN)}
\begin{algorithmic}[1]
\STATE \textbf{Input:} initial point $x_0\in\mathbb{R}^d$, step size $\alpha\in(0,1]$, batch size $b$, damping $\lambda>0$, maximum iterations $T$.
\FOR{$t=0,1,\dots,T-1$}
    \STATE Sample a mini‑batch $S_t\subset\{1,\dots,n\}$ uniformly with $|S_t|=b$.
    \STATE Compute stochastic gradient $g_t=\nabla f_{S_t}(x_t)$.
    \STATE Compute stochastic Hessian $H_t=\nabla^2 f_{S_t}(x_t)+\lambda\Id$.
    \STATE Solve $p_t=H_t^{-1}g_t$ (e.g. by CG or direct factorisation).
    \STATE Update $x_{t+1}=x_t-\alpha p_t$.
\ENDFOR
\STATE \textbf{Output:} $x_T$ (or the best iterate according to $f$).
\end{algorithmic}
\end{algorithm}

\section*{Dry Run Example}
Consider the one‑dimensional quadratic $f(w)=(w-3)^2$ with exact Hessian $2$.  
We set $\alpha=0.1$, $\lambda=0$ (so $H_t=2$ exactly), and run three iterations.

\begin{tabular}{c|c|c|c}
$t$ & $w_t$ & $g_t=\nabla f(w_t)$ & $w_{t+1}=w_t-\alpha H_t^{-1}g_t$ \\ \hline
0 & $0$ & $-6$ & $0-0.1\cdot(1/2)(-6)=0.3$ \\
1 & $0.3$ & $-5.4$ & $0.3-0.1\cdot(1/2)(-5.4)=0.57$ \\
2 & $0.57$ & $-4.86$ & $0.57-0.1\cdot(1/2)(-4.86)=0.813$ \\
\end{tabular}

Objective values: $f(0)=9$, $f(0.3)=7.29$, $f(0.57)=5.90$, $f(0.813)=4.78$, showing descent.

\section*{Assumptions}
\begin{assumption}[Smoothness]\label{as:smooth}
$f$ has $L$‑Lipschitz gradient:
\[
\|\nabla f(x)-\nabla f(y)\|\le L\|x-y\|,\qquad\forall x,y\in\mathbb{R}^d.
\]
\end{assumption}
\begin{assumption}[Hessian Lipschitz]\label{as:hessian}
The Hessian is $\rho$‑Lipschitz:
\[
\|\nabla^2 f(x)-\nabla^2 f(y)\|\le \rho\|x-y\|,\qquad\forall x,y.
\]
\end{assumption}
\begin{assumption}[Unbiased stochastic gradient]\label{as:grad-unbiased}
\[
\E[g_t\mid \mathcal{F}_t]=\nabla f(x_t),
\]
where $\mathcal{F}_t=\sigma\{x_0,S_0,\dots,S_{t-1}\}$.
\end{assumption}
\begin{assumption}[Bounded variance of gradient]\label{as:grad-var}
\[
\E\big[\|g_t-\nabla f(x_t)\|^2\mid\mathcal{F}_t\big]\le\sigma_g^2.
\]
\end{assumption}
\begin{assumption}[Unbiased stochastic Hessian]\label{as:hess-unbiased}
\[
\E[H_t\mid\mathcal{F}_t]=\nabla^2 f(x_t)+\lambda\Id.
\]
\end{assumption}
\begin{assumption}[Bounded second‑moment of Hessian inverse]\label{as:hess-inv}
There exists $c_H>0$ such that
\[
\E\big[\|H_t^{-1}\|^2\mid\mathcal{F}_t\big]\le c_H^2.
\]
\end{assumption}
All constants $L,\rho,\sigma_g,c_H,\lambda$ are assumed known and finite.

\section*{Main Convergence Theorem}
\begin{theorem}[Expected Stationarity]\label{thm:main}
Assume \ref{as:smooth}–\ref{as:hess-inv} hold and the step size satisfies $0<\alpha\le \frac{1}{c_H L}$.  
Let $x_T$ be the iterate after $T$ steps of SSN. Then  

\[
\frac1T\sum_{t=0}^{T-1}\E\big[\|\nabla f(x_t)\|^2\big]
\;\le\;
\frac{2\big(f(x_0)-f^\star\big)}{\alpha c_H T}
\;+\;
\frac{2\alpha L\sigma_g^2 c_H}{\; }.
\tag{1}
\]

Consequently, choosing $\alpha=\Theta\!\big(T^{-1/2}\big)$ yields  

\[
\min_{0\le t<T}\E\big[\|\nabla f(x_t)\|^2\big]=\mathcal{O}\!\big(T^{-1/2}\big).
\]
\end{theorem}

\section*{Theoretical Properties}
\begin{itemize}
\item \textbf{Stability.} The damping $\lambda>0$ together with bounded $\|H_t^{-1}\|$ guarantees that the Newton direction never explodes, even when the subsampled Hessian is rank‑deficient.
\item \textbf{Hyper‑parameter Sensitivity.} The bound (1) shows a linear dependence on $\alpha$ in the variance term and an inverse dependence on $\alpha$ in the bias term, leading to the optimal $\alpha\propto T^{-1/2}$.
\item \textbf{Comparison to SGD.} For the same variance $\sigma_g^2$, SSN improves the bias term by a factor $c_H\le 1$ (since $c_H$ is the expected norm of the inverse Hessian, typically $<1$ for strongly curved directions), yielding faster decay of the gradient norm.
\item \textbf{Special Cases.}
    \begin{enumerate}
        \item If $b=n$ (full batch) then $H_t=\nabla^2 f(x_t)+\lambda\Id$ and $c_H\le 1/(\lambda_{\min}+\lambda)$. The rate becomes $\mathcal{O}(T^{-1})$ under additional strong‑convexity.
        \item If $\lambda\to0$ and $H_t$ is exact, the method reduces to the classic Newton method with quadratic convergence locally.
    \end{enumerate}
\end{itemize}

\section*{Discussion}
The theorem establishes that SSN attains the same $\mathcal{O}(T^{-1/2})$ rate as stochastic gradient methods for non‑convex smooth objectives, while exploiting curvature information to reduce the constant factor. The analysis hinges on the smoothness of $f$, unbiasedness of the stochastic Hessian, and a uniform bound on the second moment of its inverse. Extensions to adaptive batch sizes or variance‑reduced Hessian estimators follow the same template.

\newpage
\section*{Preliminaries (Definitions and Lemmas)}
\begin{lemma}[Smoothness inequality]\label{lem:smooth}
If $f$ satisfies Assumption~\ref{as:smooth}, then for any $x,y$,
\[
f(y)\le f(x)+\ip{\nabla f(x)}{y-x}+\frac{L}{2}\|y-x\|^2.
\tag{2}
\]
\end{lemma}
\begin{proof}
Standard consequence of $L$‑Lipschitz gradient (integrate the gradient bound). \end{proof}

\begin{lemma}[Quadratic bound with stochastic Hessian]\label{lem:hess}
For any symmetric positive definite $H$ and vectors $g$, $x$, define $p=H^{-1}g$. Then
\[
\|p\|^2\le \|H^{-1}\|^2\|g\|^2.
\tag{3}
\]
\end{lemma}
\begin{proof}
Immediate from sub‑multiplicativity of the operator norm. \end{proof}

\section*{Main Proof (Step‑by‑Step Derivation)}
We prove Theorem~\ref{thm:main}. Throughout let $\mathcal{F}_t$ be the $\sigma$‑algebra generated by $\{x_0,S_0,\dots,S_{t-1}\}$.

\subsection*{Descent decomposition}
\begin{enumerate}[label={[\arabic*]}]
\item From Lemma~\ref{lem:smooth} with $y=x_{t+1}=x_t-\alpha H_t^{-1}g_t$,
\[
f(x_{t+1})\le f(x_t)-\alpha\ip{\nabla f(x_t)}{H_t^{-1}g_t}
       +\frac{L\alpha^2}{2}\|H_t^{-1}g_t\|^2.
\tag{4}
\]
\item Take conditional expectation w.r.t. $\mathcal{F}_t$ and use unbiasedness of $g_t$ (Assumption~\ref{as:grad-unbiased}):
\[
\E\big[f(x_{t+1})\mid\mathcal{F}_t\big]
\le f(x_t)-\alpha\ip{\nabla f(x_t)}{ \E[H_t^{-1}g_t\mid\mathcal{F}_t] }
   +\frac{L\alpha^2}{2}\E\!\big[\|H_t^{-1}g_t\|^2\mid\mathcal{F}_t\big].
\tag{5}
\]
\item Since $H_t$ and $g_t$ are conditionally independent only through $x_t$, we write
\[
\E[H_t^{-1}g_t\mid\mathcal{F}_t]
   =\E[H_t^{-1}\mid\mathcal{F}_t]\;\nabla f(x_t),
\]
because $\E[g_t\mid\mathcal{F}_t]=\nabla f(x_t)$. Define 
\[
\bar H_t^{-1}:=\E[H_t^{-1}\mid\mathcal{F}_t].
\tag{6}
\]
Thus (5) becomes
\[
\E\big[f(x_{t+1})\mid\mathcal{F}_t\big]
\le f(x_t)-\alpha\ip{\nabla f(x_t)}{\bar H_t^{-1}\nabla f(x_t)}
   +\frac{L\alpha^2}{2}\E\!\big[\|H_t^{-1}g_t\|^2\mid\mathcal{F}_t\big].
\tag{7}
\]
\item By positive definiteness of $H_t$, $\bar H_t^{-1}\succeq 0$, hence
\[
\ip{\nabla f(x_t)}{\bar H_t^{-1}\nabla f(x_t)}
   \ge \lambda_{\min}(\bar H_t^{-1})\|\nabla f(x_t)\|^2.
\tag{8}
\]
Using the bound $\|H_t^{-1}\|\le c_H$ a.s. (Assumption~\ref{as:hess-inv}) implies $\lambda_{\min}(\bar H_t^{-1})\ge 1/c_H$.
Therefore,
\[
\ip{\nabla f(x_t)}{\bar H_t^{-1}\nabla f(x_t)}
   \ge \frac{1}{c_H}\|\nabla f(x_t)\|^2.
\tag{9}
\]
\item For the variance term, expand $g_t=\nabla f(x_t)+\xi_t$ with $\E[\xi_t\mid\mathcal{F}_t]=0$ and $\E[\|\xi_t\|^2\mid\mathcal{F}_t]\le\sigma_g^2$ (Assumption~\ref{as:grad-var}). Then
\[
\|H_t^{-1}g_t\|^2
   =\|H_t^{-1}\nabla f(x_t)\|^2
     +2\ip{H_t^{-1}\nabla f(x_t)}{H_t^{-1}\xi_t}
     +\|H_t^{-1}\xi_t\|^2.
\tag{10}
\]
Taking conditional expectation, the cross term vanishes, and using Lemma~\ref{lem:hess} together with $\|H_t^{-1}\|\le c_H$,
\[
\E\!\big[\|H_t^{-1}g_t\|^2\mid\mathcal{F}_t\big]
   \le c_H^2\|\nabla f(x_t)\|^2 + c_H^2\sigma_g^2.
\tag{11}
\]
\item Substitute (9) and (11) into (7):
\[
\E\big[f(x_{t+1})\mid\mathcal{F}_t\big]
\le f(x_t)-\frac{\alpha}{c_H}\|\nabla f(x_t)\|^2
   +\frac{L\alpha^2c_H^2}{2}\big(\|\nabla f(x_t)\|^2+\sigma_g^2\big).
\tag{12}
\]
\item Collect the gradient‑norm terms:
\[
\E\big[f(x_{t+1})\mid\mathcal{F}_t\big]
\le f(x_t)-\Big(\frac{\alpha}{c_H}-\frac{L\alpha^2c_H^2}{2}\Big)\|\nabla f(x_t)\|^2
   +\frac{L\alpha^2c_H^2}{2}\sigma_g^2.
\tag{13}
\]
\item Choose $\alpha$ such that $\frac{\alpha}{c_H}-\frac{L\alpha^2c_H^2}{2}\ge \frac{\alpha}{2c_H}$, i.e.
\[
\alpha\le\frac{1}{c_H L}.
\tag{14}
\]
Under this condition,
\[
\E\big[f(x_{t+1})\mid\mathcal{F}_t\big]
\le f(x_t)-\frac{\alpha}{2c_H}\|\nabla f(x_t)\|^2
   +\frac{L\alpha^2c_H^2}{2}\sigma_g^2.
\tag{15}
\]
\item Take total expectation and sum over $t=0,\dots,T-1$:
\[
\E[f(x_T)]\le f(x_0)-\frac{\alpha}{2c_H}\sum_{t=0}^{T-1}\E\big[\|\nabla f(x_t)\|^2\big]
   +\frac{L\alpha^2c_H^2}{2}\sigma_g^2 T.
\tag{16}
\]
Rearrange:
\[
\frac{1}{T}\sum_{t=0}^{T-1}\E\big[\|\nabla f(x_t)\|^2\big]
\le \frac{2c_H\big(f(x_0)-\E[f(x_T)]\big)}{\alpha T}
   +L\alpha c_H^2\sigma_g^2.
\tag{17}
\]
Since $f(x_T)\ge f^\star$, we obtain the bound (1) of Theorem~\ref{thm:main}.
\end{enumerate}

\subsection*{Rate Derivation (With Constants)}
From (17) we have
\[
\frac{1}{T}\sum_{t=0}^{T-1}\E\big[\|\nabla f(x_t)\|^2\big]
\le \frac{2c_H\Delta_0}{\alpha T}+L\alpha c_H^2\sigma_g^2,
\quad\Delta_0:=f(x_0)-f^\star.
\]
Setting $\alpha=\sqrt{\frac{2\Delta_0}{L c_H\sigma_g^2}}\;T^{-1/2}$ (which respects (14) for sufficiently large $T$) yields
\[
\frac{1}{T}\sum_{t=0}^{T-1}\E\big[\|\nabla f(x_t)\|^2\big]
\le 2\sqrt{2L\Delta_0\sigma_g^2}\;c_H^{3/2}\;T^{-1/2}.
\]
Thus $\min_{t<T}\E[\|\nabla f(x_t)\|^2]=\mathcal{O}(T^{-1/2})$.

\section*{Special Cases}
\begin{enumerate}[label={\arabic*.}]
\item \textbf{Full‑batch Newton.} If $b=n$ and $\lambda=0$, then $H_t=\nabla^2 f(x_t)$. Under strong convexity $\mu I\preceq\nabla^2 f(x)\preceq L I$, we have $c_H=1/\mu$ and the deterministic Newton step enjoys a linear rate locally; the stochastic bound reduces to the classic deterministic descent inequality.
\item \textbf{Large damping.} If $\lambda\gg L$, then $H_t\approx\lambda\Id$, $c_H\approx 1/\lambda$, and SSN behaves like SGD with step size $\alpha/\lambda$, recovering the standard $\mathcal{O}(T^{-1/2})$ rate.
\item \textbf{Increasing batch size.} If $b$ grows with $t$ such that $\sigma_g^2=O(1/b)$, the variance term in (17) vanishes, and the optimal $\alpha$ can be taken constant, leading to an $\mathcal{O}(1/T)$ decay of the gradient norm.
\end{enumerate}

\end{document}